\begin{solapartat}
\[ \vec{F} = \vec{\mathcal{E}}\,q = m\,\vec{a}\ \text{\punts{0,5}} \Rightarrow \vec{\mathcal{E}} = \frac{m\vec{a}}{q} = \frac{9,11\times 10^{-31}\,\mathrm{kg}\cdot 1,20\times 10^{13}\,\mathrm{m/s^2}\,\vec{\imath}}{-1,60\times 10^{-19}\,\mathrm{C}} \]
\[ = -68,3\,\vec{\imath}\ \mathrm{N/C}\ \text{ó}\ \mathrm{V/m}\ \text{\punts{0,5}} \]
\end{solapartat}

\begin{solapartat}
Al ser el camp elèctric constant:
\[ \Delta V = -\vec{\mathcal{E}}\,\Delta\vec{r}\ \text{\punts{0,2}} = -(-68,3\,\vec{\imath})(0,3\,\vec{\imath}) = 20,5\ \mathrm{V}\ \text{\punts{0,2}} \]

El potencial mes alt serà a la part dreta de la càmera \punts{0,2}
\[ \Delta E = \Omega = -\Delta V\,q = -20,5\,(-1,60\times 10^{-19}) = 3,28\times 10^{-18}\ \mathrm{J} = 20,5\ \mathrm{eV}\ \text{\punts{0,4}} \]
\end{solapartat}
